Eieren in manden.
Ovens die fel branden.

Bloem in jute zakken.
Boter voor het bakken.

Muren bedekt met meel.
Potten vol met kaneel,

kruidnagel, nootmuskaat
en honingderivaat

stonden in de kasten.
Net als voor de gasten

bestemde producten.
Specerijen plukten

de bakkers liever vers,
werden in de droogpers

snel gedehydrateerd
en zo geconserveerd.

Suiker lag in een ton,
waar men handig bij kon.

Met beelden in kleuren
beschilderde deuren

zwaaiden steeds aan de kant,
zodat elke passant

van de geuren genoot
van versgebakken brood

en zoete lekkernij
rondom de bakkerij.

Ook door de raamluiken
kon men buiten ruiken

dat binnen houtgestookt
snoeperij werd gekookt.

Overal klonk rumoer.
Een leistenen werkvloer,

met cementvoeg verfraaid,
was met kruimels bezaaid

en soms stukjes wafel.
Op een grote tafel

werd deeg vooraf gekneed,
voordat het beheerst gleed

in een natte oven,
om daar klaar te stoven.

De bakker, dik en kort,
gekleed in een wit schort,

was druk aan het praten
om over te laten

zijn volleerde ambacht
aan slechts één arbeidskracht.

Terwijl hij uitleg kreeg,
kneedde hij een stuk deeg.

Eieren van pluimvee,
dat vrij op deze stee

mocht uitlopen en gaan,
werden er bij gedaan.

Toegevoegd werd honing
als suikervervanging

en voor het aroma.
Kruiden volgden erna.

Men bracht het banketdeeg,
zodat het profiel kreeg,

in een grote bakvorm.
Zo werd de taart conform

de wensen van de klant.
In de oven beland,

werd op de tijd gelet.
Dit unieke banket

kon wel een uur duren
en moest vele uren

gradueel afkoelen.
De bakker kon voelen

of de taart gereed was.
Dan verschoof hij het pas.

Darmintelligentie
als zelfregulatie

was minder bij Greta.
Overeten kwam na

lastige emoties
als er geen restricties

op haar waren gelegd.
Er werd zelfs soms gezegd

dat ze werd aangespoord
tot gulzigheidzelfmoord.

Haar serotonine
als geluksbenzine

vulden haar darmen aan
door voor vraatzucht te gaan.

Daar ze nooit genoeg kreeg
bleven haar darmen leeg.

Zo ontving Greta’s brein
van haar pens een noodsein,

omdat die onbewust
haar ooit zo had gesust.

Uit schaamte vroeg gespeend
werd Greta snel getraind

om haar zuigreflexen
gemaakt te relaxen

via haar rechterduim.
Moeder zag haar verzuim

echter niet als reden.
Dus werd er gestreden

tegen het duimzuigen.
Greta moest wel buigen.

Geestelijk verwaarloosd
werd haar enige troost

overmatig eten
om stress te vergeten.

Haar impulsiviteit
bracht Greta daarna spijt.

bij gebrek aan troosten
van hen die

steeds door.

Haar darmen bleven leeg
omdat

Als

Zeker, het concept van darmintelligentie (ofwel de "gut brain connection") kan een belangrijke rol spelen in Greta's vraatzucht en overeten. Darmintelligentie verwijst naar het complexe netwerk van zenuwen, hormonen en bacteriën in de darm die signalen naar de hersenen sturen en omgekeerd, vaak aangeduid als de hersenen van de darmen. Als deze verbinding uit balans is, kan dat leiden tot emotioneel eten en een gebrekkige zelfregulatie.

Voor Greta zou een gebrek aan darmintelligentie kunnen betekenen dat:

Ze niet goed voelt wat haar lichaam nodig heeft: Greta kan de natuurlijke signalen van honger en verzadiging niet herkennen of erop reageren. Dit kan het gevolg zijn van een verstoring in haar darm-hersencommunicatie door ondervoeding, stress, of trauma in haar jeugd, zoals de harde methoden waarmee ze werd opgevoed.

Ze eten gebruikt om emoties te reguleren: Een gebrekkige darmintelligentie kan leiden tot een overreactie op stress of negatieve emoties. De darmen produceren immers een groot deel van het lichaamseigen serotonine (het "gelukshormoon"). Greta kan eten gebruiken om dit tekort te compenseren, waardoor ze een cyclus van troost en schuld ontwikkelt.

Haar darmen symbolisch leeg zijn: Greta's vraatzucht kan ook symbool staan voor een innerlijke leegte of gebrek aan troost in haar vroege leven. Het snelle spenen en de sociale druk van haar moeder kunnen ervoor gezorgd hebben dat ze nooit
